\documentclass[11pt]{article}
\usepackage[T1]{fontenc}
\usepackage[utf8]{inputenc}
\usepackage{lmodern,microtype}
\usepackage[a4paper,margin=25mm]{geometry}
\usepackage{amsmath,amssymb,amsthm,mathtools,mathrsfs}
\usepackage{booktabs,longtable,array,enumitem}
\usepackage{xcolor,hyperref,fancyhdr}
\hypersetup{colorlinks=true,linkcolor=blue!45!black,urlcolor=blue!45!black,citecolor=blue!45!black,pdftitle={Fifteen integer unknowns and smaller Diophantine forms for Con(ZFC)}}
\setlength{\headheight}{14pt}
\pagestyle{fancy}\fancyhf{}\lhead{Diophantine forms of consistency}\rhead{ProveIt synthesis}\cfoot{\thepage}
\newtheorem{theorem}{Theorem}[section]
\newtheorem{lemma}[theorem]{Lemma}
\newtheorem{proposition}[theorem]{Proposition}
\newtheorem{corollary}[theorem]{Corollary}
\theoremstyle{definition}\newtheorem{definition}[theorem]{Definition}
\newcommand{\Z}{\mathbb Z}\newcommand{\N}{\mathbb N}\newcommand{\Q}{\mathbb Q}
\newcommand{\Sq}{\operatorname{Sq}(\Z)}
\newcommand{\Con}{\operatorname{Con}(\mathrm{ZFC})}
\newcommand{\pin}{\texttt{8c6eec9d19ef80538c8e61ff6e0180040583446a}}
\newcommand{\smallpath}[1]{{\small\path{#1}}}
\allowdisplaybreaks
\newcommand{\OpsSixteen}{1,550}
\newcommand{\OpsFifteen}{1,541}
\newcommand{\DegreeFifteenFormula}{\ensuremath{813056\,5^{33}+4184}}

\title{Fifteen integer unknowns and smaller coefficient tuples\\for Diophantine forms of ZFC consistency}
\author{AI-assisted research synthesis for the ProveIt repository}
\date{22 October 2026}
\begin{document}
\maketitle
\begin{abstract}
We compare the seven current constructions in ProveIt's collection and combine their ideas in three ways. First, a Pell orbit argument removes an unnecessary sign guard from the signed Jones construction, giving an explicit polynomial in sixteen unrestricted integers. Restricting the coding radix to a cofinal sequence of square multiples of a fixed power of two removes a second guard, giving a polynomial in fifteen unrestricted integers. Both are delivered as division-free arithmetic circuits with completely specified fixed coefficient recipes. A smaller separating weight and a factored evaluator make the sixteen-variable circuit substantially cheaper than its seventeen-variable predecessor. Second, three exact identities reduce the existing degree-1081 twelve-positive-variable evaluator from 111 to 108 operations; a transport shift gives a 103-operation three-equation system. Third, combining the mixed-field layout with arbitrary-factor Boolean padding reduces the compiler base from $2^{1953125}$ to $2^{1064921}$ and materializes all four coefficient ports in 533,099 uncompressed bytes. Each construction has the same solvability condition as the calibrated blank K215 run. Thus ZFC proves its unsolvability equivalent to $\Con$. This is an ordinary mathematical reduction with executable finite audits; no new end-to-end proof-assistant certificate or global optimality claim is made.
\end{abstract}
\setcounter{tocdepth}{1}
\tableofcontents
\newpage

\section{Outcome, scope, and comparison}
\subsection{What the statement means}
Throughout, $\N=\{0,1,2,\ldots\}$ and ``positive'' means strictly positive. The principal new result is a particular polynomial $\mathcal P_{15}\in\Z[x_1,\ldots,x_{15}]$, specified in Sections~\ref{sec:base}--\ref{sec:fifteen} and in the delivered circuit. Its theorem is
\begin{equation}\label{eq:main}
 \mathrm{ZFC}\vdash
 \left(\neg\exists\mathbf x\in\Z^{15}:\mathcal P_{15}(\mathbf x)=0\right)
 \longleftrightarrow\Con.
\end{equation}
Here $\Con$ is the usual arithmetical assertion that there is no finite ZFC refutation in the proof system used by the repository's calibration. We establish an equivalence; neither side is assumed or decided. In particular, equation~\eqref{eq:main} is compatible with the second incompleteness theorem. We make no two-sided independence claim from consistency alone.

The reduction has three levels. The literal K215 machine is calibrated to finite ZFC refutations; the native compiler turns blank halting into positive Diophantine solvability; the new algebraic and number-theoretic transformations reduce the complexity of that presentation. The first two levels are inherited ordinary mathematical results documented in the repository. Its Lean and Rocq endpoint work provides valuable provenance, but the native compiler and the new transformations have not been joined into a newly kernel-checked theorem. The finite programs in this archive audit formulas, coefficients, counts, and selected identities. They do not prove an unbounded simulation theorem by testing examples.

\subsection{A fixed comparison snapshot}
All comparisons refer to ProveIt commit \pin. The collection at that snapshot contains packages 24--30, with its own reconciliation note and provenance warnings~\cite{collection}. It is essential to distinguish the number of unknowns, their domain, degree, arithmetic evaluation cost, and coefficient height. None is an adequate substitute for all the others.

\begin{center}\small
\begin{tabular}{p{51mm}rrp{43mm}}
\toprule
Construction & Unknowns & Operations & Degree or qualification\\
\midrule
Package 29 signed Jones & $17\Z$ & 20,607 & $424296\,5^{33}+53824$\\
New signed Jones & $16\Z$ & \OpsSixteen & $406528\,5^{33}+10008$\\
New square-radix Jones & $15\Z$ & \OpsFifteen & \DegreeFifteenFormula\\
Package 30 positive product & 12 positive & 111 & 1081\\
New positive product evaluator & 12 positive & 108 & 1081\\
Package 30 positive system & 12 positive & 107 & 3 equations, maximum 1072\\
New shifted positive system & 12 positive & 103 & 3 equations, maximum 1072\\
\bottomrule
\end{tabular}
\end{center}
Operations mean live binary additions, subtractions, and multiplications, including multiplication by a fixed integer. Reuse and equality comparisons are free. Fixed coefficient construction is outside this arithmetic evaluation count, as in the source packages. These counts are explicit evaluator sizes, not circuit lower bounds or bit-complexity estimates.

Several existing results retain different advantages. Package 24 has twelve positive unknowns at degree 1069. Package 30 has a 31-positive-variable quartic and a nine-natural-variable specialization of degree $120724\,5^{32}+366$. Package 27 has a nine-positive-variable circuit and low-degree systems using 59 operations. The inherited fifteen-positive-variable $P_{67}$ remains the collection's 67-operation single equation. The new fifteen-integer result does not reduce the existing nine-positive-variable count; the domains differ. It also does not improve the best low-degree forms.

The literature has a smaller general integer-variable bound: Sun's September 2026 preprint states a ten-integer representation theorem for recursively enumerable sets~\cite{sun}. We verified the source statement, but do not instantiate its different compiler here. Consequently ``fifteen'' is an improvement over the inspected concrete repository construction, not a universal record or a proof of global minimality.

\subsection{How the ideas are combined}
The integer construction uses package 29's exact T4-bound Jones body, its coprime Pell guards, and its degree-five intersective divisor merger. It adds a sign-normalization theorem for the final Pell coordinate, a cofinal square-radix specialization, and a smaller separating weight. Its efficient evaluation uses the algebraic digit-polynomial organization developed in package 27, with the original T4 coefficients retained. The positive construction uses package 30's polynomial and proves three identities in the polynomial ring before binding all four ports to a complete compiler tuple. The smaller tuple combines package 25's mixed-field layout with package 24's arbitrary-factor completeness controller. Every resulting family has a separately named coefficient binding.

\section{The fixed Jones body}\label{sec:base}
\subsection{A finite recipe for every coefficient}
For the integer constructions we retain exactly the T4 instance used by package 29's signed17 construction~\cite{signed17}. Its source quartic has $\nu=32$ positive coordinates. Shift each by one to obtain a polynomial $P$ over $\N^{32}$; add a dummy input coordinate, fixed to one. The source proves the normalization $P(1,0,\ldots,0)\ne0$ and
\[
 \exists\mathbf n\in\N^{32}:P(1,\mathbf n)=0
 \quad\longleftrightarrow\quad\text{K215 halts on its blank tape}.
\]
Let $p_\alpha$ be the coefficient of $P$ at exponent vector $\alpha\in\N^{33}$, $|\alpha|\le4$, and define
\[
 c_\alpha=(4-|\alpha|)!\prod_{r=0}^{32}\alpha_r!\,p_\alpha,
 \qquad L=5^{33}.
\]
The supplied source data give explicit polynomial bounds $\rho_j$ for the shifted quadratic residuals. Use the inherited fixed recipe
\begin{align*}
 s_*&=2+24\sum_j\rho_j^2,& z_*&=2^{s_*},\\
 e_*(T)&=\sum_{|\alpha|\le4}(z_*+c_\alpha)
       T^{L-\sum_{r=0}^{32}\alpha_r5^r},& y_*&=e_*(2z_*).
\end{align*}
These are fixed integers and a fixed polynomial, not additional unknowns. Terms with $c_\alpha=0$ still occur through the $z_*$ summand. The bound proves $|c_\alpha|+1<z_*$ for every coefficient. All exponents are fixed nonnegative integers.

The archive contains the literal quartic, the coefficient polynomials, the complete T4 recipe, the ordered transition/window data, the source licences, and hashes authenticating the imported files. Thus these symbols do not hide an unspecified index or a search for a halting computation. The T4 and Jones integers are specified by finite recipes rather than expanded into infeasibly large decimal strings. The materialized small tuple in Section~\ref{sec:coeff} belongs to the positive family; it is not silently substituted into this integer construction.

\subsection{Eight base coordinates and their expressions}
For the sixteen-integer construction take
\[
 (\varepsilon,g_0,h_0,s_0,w_0,\phi,i_0,j_0)\in\Z^8,
\]
and put $g=g_0+1$, $h=h_0+1$, $s=s_0+1$, $w=w_0+1$, $i=i_0+1$, $j=j_0+1$. Define
\begin{gather*}
 b=y_*+1+\varepsilon,\quad \mathcal B=b^5,\quad q=\mathcal B^L,
 \quad\ell=\sum_{r=1}^{32}\mathcal B^{5^r},
 \quad\lambda=\sum_{r=0}^{4L-1}\mathcal B^r,\\
 e=e_*(\mathcal B),\quad\theta=\mathcal B-2z_*,\quad
 N=q^{16},\quad G=32b\mathcal B^{5^{32}}.
\end{gather*}
The two-bracket formula for the packed integer is
\begin{align}
 R={}&\bigl[g+eq^3+\ell q^5+
 \bigl(2(e-z_*\lambda)(1+\mathcal B+g)^4
       +\lambda\mathcal B+\lambda\mathcal Bq^4\bigr)q^7\bigr](N^2-N)
 \nonumber\\
 &+\bigl[q^3-b\ell+\ell+\theta\lambda q^3+
             (\mathcal B-2)q^8\bigr](N^2-1).\label{eq:R}
\end{align}
The remaining expressions are
\begin{gather*}
 U=N^2w,\quad Y=N^2s,\quad M=RY,\quad P_0=2M^2U,\\
 K=R+1+h(P_0-1),\quad A=M(U+1),\quad B_1=2R+1,\quad W_1=bw,\\
 C=B_1+W_1+\phi,\quad J=4A-5,\quad Q=3W_1C-2(W_1^2-1),\\
 D=(A^2-1)C^2+1,\quad E=iJDC^2,\quad F=(A^2-1)E^2+1,\\
 G_1=A+F(CD+1-A),\quad H=B_1+jC,\quad I=(G_1^2-1)H^2+1,\\
 a_0=(P_0^2-1)K^2+1,\quad a_1=DFI,\\
 m=K^2G-4(C-KY)^2G-8gK^2,\quad D_0=JF,\quad T_0=QF+J(H-C).
\end{gather*}
Every intermediate is an abbreviation, not a quantified variable. The source ratio theorem states that the source quartic has a solution exactly when eight nonnegative base coordinates satisfy
\begin{equation}\label{eq:JR}
 a_0,a_1\in\Sq,\qquad D_0\mid T_0,\qquad m>0.
 \tag{JR}
\end{equation}
The repository's parametric nine-variable work records this route through \texttt{NineReduce} and \texttt{NineRatio}~\cite{ninelean}. We use this specific interface, not a claim that every nine-variable polynomial automatically has the required extra structure.

\section{Removing the sign guard on the final Pell coordinate}\label{sec:pell}
\subsection{The bounds used before normalization}
Only the five coordinates
\[
 \varepsilon,\ g_0,\ s_0,\ w_0,\ \phi
\]
will initially be required nonnegative. In addition impose $i\ne0$, $m>0$, and the square/divisibility conditions in~\eqref{eq:JR}. The inherited signed-$h$ argument remains valid with no assumption on $j$; for clarity we give its essential steps.

Since $g>0$ and $G>0$, the margin yields
\[
 K\ne0,\qquad g<G/8=4b\mathcal B^{5^{32}},
 \qquad 4(C-KY)^2<K^2.
\]
Write $t=5^{32}$. The index bounds imply $b>y_*\ge2z_*$ and $\mathcal B\ge32$. Hence
\[
 1+\mathcal B+g<\mathcal B^{t+2}\le q,
 \qquad b\ell<q^3.
\]
In~\eqref{eq:R}, the potentially negative term
$-2z_*\lambda(1+\mathcal B+g)^4$ is outweighed by
$\lambda\mathcal Bq^4$. Both brackets are positive, so $R\ge N^2-N\ge2$.
Consequently $Y,U\ge1$, $C=2R+1+bw+\phi>0$, and $P_0\ge2R^2>R+2$.
The margin forces $K>0$, since $K<0$ would give $|C-KY|>|K|$.

If $h\le0$, then $1\le K\le R+1<P_0$. The case $K=1$ would force
$P_0-1\mid R$, which is impossible. For $2\le K<P_0$,
\[
 (P_0K-1)^2<(P_0^2-1)K^2+1<(P_0K)^2,
\]
contradicting squarehood. Thus $h\ge1$. Replacing $i$ by $|i|$ preserves the five expressions in~\eqref{eq:JR}, because $i$ occurs there only through $E^2$. This gives the original nonnegative $h_0$ and $i_0$ conventions.

\begin{lemma}[Nonzero divisor for every integer assignment]\label{lem:denom}
The polynomial $D_0=JF$ never vanishes on integer base coordinates.
\end{lemma}
\begin{proof}
$J=4A-5$ cannot be zero. If $F=0$, then $(A^2-1)E^2=-1$, forcing $A=0$ and $E=\pm1$. But $J=-5$ and $E=iJDC^2$ is divisible by five. This is impossible. The argument does not presuppose any guard or square condition.
\end{proof}

\begin{lemma}[Coprimality of the three Pell radicands]\label{lem:coprime}
Under the five nonnegative guards and~\eqref{eq:JR}, with $i\ne0$, the integers $D,F,I$ are positive and pairwise coprime. In particular each is a square.
\end{lemma}
\begin{proof}
The preceding bounds give $A\ge4$, $C>0$, and $D,F>1$.
Also $CD+1-A>0$, so $G_1>1$ and $I\ge1$. Since $E$ is divisible by $D$ and $J$, $F\equiv1\pmod D$ and $F\equiv1\pmod J$. Thus $D\perp F$ and $J\perp F$.

The definition gives $G_1\equiv1\pmod D$, whence $I\equiv1\pmod D$. From $JF\mid QF+J(H-C)$ and $J\perp F$ we obtain $H\equiv C\pmod F$. As $G_1\equiv A\pmod F$, it follows that
\[
 I\equiv(A^2-1)C^2+1=D\pmod F.
\]
Therefore $I$ is coprime to $D$ and to $F$. Unique factorization and the hypothesis $DFI\in\Sq$ show that all three positive factors are squares.
\end{proof}

\begin{lemma}[A positive Pell ordinate in a prescribed residue class]\label{lem:orbit}
Let $t\ge2$ and $M_0\ge1$ be integers. If an integer $H_0$ satisfies
$r_0^2-(t^2-1)H_0^2=1$ for some integer $r_0$, then there are arbitrarily large positive integers $H'$ congruent to $H_0$ modulo $M_0$ for which $(t^2-1)H'^2+1$ is a square.
\end{lemma}
\begin{proof}
Put $\delta=t^2-1$ and choose $r_0>0$. The norm-one real number
$\alpha=r_0+H_0\sqrt\delta$ is positive, even if $H_0\le0$.
The integer matrix
\[
 U=\begin{pmatrix}t&\delta\\1&t\end{pmatrix}
\]
has determinant one, so its reduction belongs to the finite group
$\mathrm{SL}_2(\Z/M_0\Z)$. Some positive power $U^k$ is the identity modulo $M_0$. The coordinates of
\[
 \alpha(t+\sqrt\delta)^{kn}=r_n+H_n\sqrt\delta
\]
are integers, have norm one, and are congruent to $(r_0,H_0)$ modulo $M_0$.
As $n\to\infty$, this real number tends to infinity while its conjugate tends to zero. Thus $H_n\to+\infty$. This also handles $H_0=0$; no assertion about a fundamental Pell unit is needed.
\end{proof}

\begin{proposition}[The $j_0$ guard is unnecessary]\label{prop:jguard}
The existential condition~\eqref{eq:JR} is unchanged if $h_0,i_0,j_0$ are allowed arbitrary integers, provided $i\ne0$ and the other five base coordinates are nonnegative.
\end{proposition}
\begin{proof}
Normalize $h$ and $i$ as above. By Lemma~\ref{lem:coprime}, $I$ is a square. Apply Lemma~\ref{lem:orbit} with $t=G_1$, modulus $CF$, and ordinate $H$. Choose $H'>B_1+C$ and $H'\equiv H\pmod{CF}$.
Since $H=B_1+jC$, the number $j'=(H'-B_1)/C$ is an integer and $j'>1$. Replacing $j$ by $j'$ replaces $I$ by a square $I'$. It preserves $DFI\in\Sq$, and changes $T_0$ by $J(H'-H)$, which is divisible by $JF$. The quantities $a_0$ and $m$ are independent of $j$ and remain unchanged. All eight base coordinates now satisfy the original nonnegative conventions. The reverse implication is immediate.
\end{proof}

\section{One explicit polynomial in sixteen integers}\label{sec:sixteen}
\subsection{Combining inequalities by coprime Pell guards}
Let $(x_1,\ldots,x_5)=(\varepsilon,g_0,s_0,w_0,\phi)$ and introduce integer variables $z_1,\ldots,z_6$. Set $C_0=1$ and
\begin{align}
 C_r&=C_{r-1}\bigl(1+4(3x_r+2)C_{r-1}^2z_r^2\bigr)
       &&(1\le r\le5),\nonumber\\
 C_*&=C_5\bigl(1+4(3m-1)C_5^2z_6^2\bigr).\label{eq:guards16}
\end{align}
Subject to $z_1\cdots z_6\ne0$, the assertion $C_*\in\Sq$ is equivalent to the existence of the five nonnegative guards and $m>0$, in the following precise sense: every such square forces the inequalities; if the inequalities hold, nonzero $z_r$ can be chosen to make the product a square.

Indeed, each factor is one modulo four and coprime to every preceding factor. Squarehood of the product makes the absolute value of every factor a square; a negative square cannot be one modulo four, so every factor is a positive square. A negative discriminant $3x_r+2$ or $3m-1$ would make its factor negative. These discriminants are two modulo three and cannot be zero. Thus $x_r\ge0$ and $m\ge1$.
Conversely, each positive discriminant is two modulo three, hence nonsquare. Its Pell equation has a nontrivial solution with ordinate divisible by any prescribed positive integer, by the same finite-matrix argument as Lemma~\ref{lem:orbit}. Choose it divisible by $2|C_{r-1}|$ and set $z_r$ to the quotient. Iteration makes every factor a square.

\subsection{An intersective divisor merger}
Use the fixed polynomial
\[
 \Pi(v)=(v^3-v-1)(v^2+23).
\]
It has no rational root but a root modulo every positive integer. Here is a complete short proof. The rational root test excludes roots of the cubic, and the quadratic has no real root. For an odd prime $p\ne23$, either the cubic has a root, necessarily simple, or it is irreducible. In the latter case Frobenius acts on its roots as a three-cycle, so the product of their pairwise differences is in $\mathbb F_p$; its square is the discriminant $-23$. Thus the quadratic has a simple root. Simple roots lift to prime powers. At $p=23$, the cubic has the simple root $3$. At $p=2$, $-23\equiv1\pmod8$, and the usual odd-square lifting gives roots of the quadratic modulo every power of two. The Chinese remainder theorem completes the proof.

For $D,T,v\in\Z$ define the homogeneous polynomial
\begin{equation}\label{eq:Rhom}
 \mathscr R(D,T,v)=
 \bigl((T+Dv)^3-D^2(T+Dv)-D^3\bigr)
 \bigl((T+Dv)^2+23D^2\bigr).
\end{equation}
If $D\ne0$, it never vanishes, since it is $D^5\Pi(T/D+v)$.
For integers $D\ne0,T,Z$,
\begin{equation}\label{eq:merge}
 D\mid T\ \land\ Z\ne0
 \quad\Longleftrightarrow\quad
 \exists v\in\Z:\ ZD^5\mid\mathscr R(D,T,v).
\end{equation}
The forward implication follows from intersectivity. Conversely, $Z=0$ would force the nonzero right-hand polynomial to vanish. If $D\nmid T$, choose a prime with $\alpha=v_p(D)>\beta=v_p(T)$. Then $v_p(T+Dv)=\beta$; the first factor in~\eqref{eq:Rhom} has unique smallest valuation $3\beta$, and the second has valuation $2\beta$. Their product has valuation $5\beta<5\alpha$, a contradiction.

\subsection{A smaller separating weight}
The guard product is always one modulo four. For every integer $C_*\equiv1\pmod4$,
\begin{equation}\label{eq:smallweight}
 |C_*+1|>\sqrt{|C_*|}.
\end{equation}
For $C_*\ge1$ this is immediate; for negative $C_*$, one has $C_*\le-3$ and $(|C_*|-1)^2>|C_*|$.
Also, for all real $P_0,K$,
\[
 \sqrt{|(P_0^2-1)K^2+1|}\le P_0^2+K^2+1.
\]
Consequently, with
\begin{equation}\label{eq:weights}
 w_1=C_*+1,\qquad w_2=w_1(P_0^2+K^2+2),
\end{equation}
we have
\[
 |w_1|>\sqrt{|C_*|},\qquad
 |w_2|>\sqrt{|C_*|}+|w_1|\sqrt{|a_0|}.
\]
The source used $w_1=1+C_*^2$. The linear choice is valid because the guard product already has its special residue class. Only magnitudes enter the separation proof, so negative weights create no problem.

\subsection{The division-free final formula}
Introduce two further integer variables $v,u$, and define
\begin{gather*}
 Z=i z_1z_2z_3z_4z_5z_6,\qquad
 S=ZD_0^5,\qquad T=\mathscr R(D_0,T_0,v),\qquad X=Su-T,\\
 A_0=S^2C_*,\qquad A_1=S^2w_1^2a_0,\qquad A_2=S^2w_2^2a_1,\\
 V=X^2+A_2-A_0-A_1,\qquad
 W=V^2+4X^2A_2-4A_0A_1.
\end{gather*}
Our polynomial is
\begin{equation}\label{eq:finalpoly}
 \boxed{\mathcal P=W^2-16V^2X^2A_2.}
\end{equation}
With~\eqref{eq:guards16} this is $\mathcal P_{16}$. It has exactly eight base coordinates, six guard variables, and $v,u$: sixteen unrestricted integers.

\begin{lemma}[The finalizer]\label{lem:finalizer}
For fixed earlier coordinates, there exists an integer $u$ making~\eqref{eq:finalpoly} zero if and only if
\[
 S\ne0,\qquad S\mid T,\qquad C_*,a_0,a_1\in\Sq.
\]
\end{lemma}
\begin{proof}
By Lemma~\ref{lem:denom}, $D_0\ne0$, and~\eqref{eq:Rhom} implies $T\ne0$. If $S=0$, then $A_0=A_1=A_2=0$, $X=-T$, and $\mathcal P=T^8\ne0$.
For square roots $r_0^2=C_*$, $r_1^2=a_0$, $r_2^2=a_1$ in a complex algebraic extension, direct multiplication gives
\[
 \mathcal P=\prod_{\epsilon_0,\epsilon_1,\epsilon_2=\pm1}
 \bigl[X-S(\epsilon_0r_0+\epsilon_1w_1r_1+\epsilon_2w_2r_2)\bigr].
\]
A zero makes one encoded radical sum rational. For two distinct sign patterns, take the last nonzero root whose sign changes. Its absolute value is at least one, and its weight strictly exceeds the sum of the absolute values of all earlier contributions, by~\eqref{eq:weights}. Thus distinct effective sign patterns give distinct sums, including when some radicands are negative.
The field generated by the roots is a finite Galois extension of $\Q$. Every automorphism acts by sign changes. Rationality of the encoded sum and separation force every root in it to be fixed, hence rational, hence an integer. Thus the three radicands are nonnegative squares, and $u-T/S$ is an integer, proving $S\mid T$. Conversely, choose integer roots and set $u=T/S+r_0+w_1r_1+w_2r_2$.
\end{proof}

\begin{theorem}[Sixteen unrestricted integers]\label{thm:16}
The polynomial $\mathcal P_{16}$ is solvable over $\Z$ exactly when the pinned source quartic is solvable. Its exact total degree is $406528\,5^{33}+10008$, and its supplied circuit uses \OpsSixteen\ operations.
\end{theorem}
\begin{proof}
A zero gives the three squares and divisibility by Lemma~\ref{lem:finalizer}. Equation~\eqref{eq:merge} gives $i\ne0$, nonzero guard witnesses, and $D_0\mid T_0$. The guard lemma gives the five inequalities and $m>0$. Proposition~\ref{prop:jguard} yields the original source predicate. Conversely, a source witness supplies the guards, nonzero Pell witnesses, an intersective residue $v$, and the final integer $u$. The complexity assertions are justified in Section~\ref{sec:verification}.
\end{proof}

\section{Removing a second guard: fifteen integers}\label{sec:fifteen}
\subsection{A cofinal family of admissible radices}
The coefficient-digit bound gives
\begin{equation}\label{eq:bfloor}
 y_*<(2z_*)^{L+1}=:b_*.
\end{equation}
Indeed the digits of $e_*$ at base $2z_*$ all lie between zero and $2z_*-1$, and its degree is at most $L$. Since $z_*$ is a power of two, $b_*$ is also a fixed power of two.

\begin{lemma}[Cofinal radix choice]\label{lem:cofinal}
Given a solution of the normalized source quartic, the Jones reduction can be completed using every sufficiently large power of two as $b$. In particular it can be completed with
\[
 b=b_*2^{2t}
\]
for a sufficiently large natural number $t$.
\end{lemma}
\begin{proof}
Let $(1,n_1,\ldots,n_{32})$ be a source solution. In the source necessity construction, choose a power of two $b$ larger than $y_*$ and all these coordinates. The code is
\[
 g=\sum_{r=1}^{32} n_r(b^5)^{5^r}.
\]
Normalization makes $g>0$, and the digit bound gives
$g<b(b^5)^{5^{32}}$. The definitions of $e,\ell,q,\lambda,\theta$ and $R$ are the same polynomial expressions as above. The transfer and carry arguments use $b=2^k$ and the stated strict lower bounds, so they yield
$N^2\mid\binom{2R}{R}$. The ratio-completeness theorem then supplies positive $h,s,w,i,j$ and nonnegative $\phi$ satisfying~\eqref{eq:JR}.

To identify the strengthening precisely, the repository's \texttt{USys\_of\_mem} chooses $k=xy_*+\sum_{r=0}^{32}n_r+1$, with $n_0=x$; at input one this is $y_*+\sum_{r=1}^{32}n_r+2$. Its subsequent proof uses that choice only to establish $k\ge1$, $y_*<b$, and each $n_r<b$. Replacing it by any larger exponent preserves those hypotheses; \texttt{reduced\_of\_mem} preserves $b$, and \texttt{ratio\_complete} accepts an arbitrary witness of $b=2^k$. No exponent congruence is imposed. The exponents of $b_*2^{2t}$ are unbounded, proving the specialization. The included source excerpts and independent review identify this exact proof dependency; a claim of unspecified cofinality would not suffice.
\end{proof}

\subsection{The fifteen-variable definition}
Replace the base coordinate $\varepsilon$ by an unrestricted integer $r$ and replace the definition of $b$ by
\begin{equation}\label{eq:squarebase}
 b=b_*r^2.
\end{equation}
Keep all other formulas in Section~\ref{sec:base}. The eight base coordinates are now
\[
 (r,g_0,h_0,s_0,w_0,\phi,i_0,j_0)\in\Z^8.
\]
Only $(x_1,x_2,x_3,x_4)=(g_0,s_0,w_0,\phi)$ need inequality guards. Introduce five integer variables $z_1,\ldots,z_5$, set $C_0=1$, and put
\begin{align*}
 C_j&=C_{j-1}\bigl(1+4(3x_j+2)C_{j-1}^2z_j^2\bigr)
            &&(1\le j\le4),\\
 C_*&=C_4\bigl(1+4(3m-1)C_4^2z_5^2\bigr),
 &Z&=ri z_1z_2z_3z_4z_5.
\end{align*}
Use the same weights~\eqref{eq:weights}, the same $S,T,X,A_0,A_1,A_2,V,W$, and the same final formula~\eqref{eq:finalpoly}. This is $\mathcal P_{15}$.
The complete variable list is
\[
 \boxed{r,g_0,h_0,s_0,w_0,\phi,i_0,j_0,z_1,z_2,z_3,z_4,z_5,v,u.}
\]
There are no quantified radical, divisibility, inequality, or exponentiation conditions in the delivered equation. Such conditions appear only in its proof. Fixed large powers are expanded by a finite addition/multiplication circuit.

\begin{theorem}[Fifteen unrestricted integers]\label{thm:15}
The equation $\mathcal P_{15}=0$ is solvable over $\Z$ if and only if the calibrated K215 machine halts on a blank tape. Therefore~\eqref{eq:main} holds. Its supplied evaluator has \OpsFifteen\ operations and exact total degree \DegreeFifteenFormula.
\end{theorem}
\begin{proof}
For soundness, the finalizer still excludes $S=0$, using the unconditional denominator lemma. Hence $r\ne0$, $i\ne0$, and all five guard variables are nonzero. Therefore $b\ge b_*>y_*$. Define
$\varepsilon=b-y_*-1\ge0$. All body expressions are now exactly those of the original signed Jones source at this nonnegative $\varepsilon$. The remaining four guards and margin follow from the stack. The signed-$h$, signed-$i$, and Pell-orbit arguments produce the original source witness, hence blank halting.

For completeness, start from blank halting and the corresponding normalized quartic witness. Lemma~\ref{lem:cofinal} gives a source ratio witness at $b=b_*2^{2t}$. Choose $r=2^t$, which is nonzero. The five positive discriminants admit nonzero Pell witnesses making $C_*$ a square. Intersectivity supplies $v$ satisfying the merged divisibility, and the radical finalizer supplies $u$. Thus the equation has an integer zero. Every argument is ordinary number theory formalizable in ZFC; composing with the calibrated refutation equivalence and negating proves the assertion.
\end{proof}

The square substitution doubles the leading degree contribution of the coding radix. The fifteen-variable result is therefore a variable-count tradeoff; the sixteen-variable companion remains the lower-degree option. The fixed integer $b_*$ has an explicit recipe and is not an extra unknown. It is not charged as a computed expression during evaluation, consistently with the source's treatment of $z_*$ and $y_*$.

\section{A smaller evaluator in twelve positive unknowns}\label{sec:positive}
For the direct positive family use package 30's twelve supplied positive variables
\[
 J,F,\alpha,z,h,i,T,s,w,\eta,\zeta,t.
\]
The following notation is local to this section. The four fixed ports are $\beta=B-1$, $K_0$, $M_C$, and the shifted $M_F$. Define
\begin{gather*}
 q=\beta J+1,\quad C=q-F-\alpha,\quad X=wq,\quad Y=sq^3,\quad E=XY,\\
 k=\eta+\zeta,\quad c=kY+\eta,\quad a=E+Y,\quad H=4a+3,\quad\Delta=a^2+H,\\
 \ell=EkY,\quad A=1+\ell(\ell+k),\quad G=1+\Delta c^2,\quad b_0=X+ac,\\
 R=((q-1)(q-F)+\alpha)(q^2-1)+(M_C+qM_F)J,\\
 N_k=k-hE-R,\quad N_t=(K_0+w)C+q-F-z(q-1).
\end{gather*}
Next put
\begin{gather*}
 I=ic^2,\quad D=1+\Delta I^2,\quad Q=\Delta(D-1),\quad d=Q-1,
 \quad v=cT,\quad b=c+RD,\\
 M=k(I^2+T^2),\quad L_0=HM+1,\quad W_0=b_0+Ht,\quad n=(HM)^2A,
 \quad d_n=dn.
\end{gather*}
Here is the optimized evaluator, with every named expression an abbreviation:
\begin{align*}
 H_Q&=H^2Q,&H_P&=H_Q(Dv^2),&H_U&=H_Qb^2,\\
 H_e&=H_P+H_U-H^2,&j&=4H_PH_U,\\
 \tau&=2d(W_0^2+L_0^2G),&\kappa&=[d(W_0^2-L_0^2G)]^2,\\
 a_a&=2(d_n+H_e),&c_0&=(d_n-H_e)^2-j,\\
 a_t&=\tau-a_a,&b_t&=c_0-\kappa,&c_f&=16d_nj,\\
 v_r&=a_t(a_t\tau+2b_t)-c_f,&u_r&=b_t^2-a_t^2\kappa,\\
 \mathcal N&=u_r(u_r+\tau v_r)+\kappa v_r^2.
\end{align*}
The equation is
\begin{equation}\label{eq:positive12}
 \boxed{N_kN_t(1+\mathcal N)-1=0.}
\end{equation}
It has degree 1081 and an explicit 108-operation circuit. The three-equation presentation $N_k=1$, $N_t=1$, $\mathcal N=0$ takes 104 operations.

\subsection{Why the saving is exact}
In the source, $H_P=H^2(QDv^2)$ and $H_U=H^2(Qb^2)$ are computed separately. Sharing $H_Q=H^2Q$ saves one multiplication. Computing $\kappa=[d(W_0^2-L_0^2G)]^2$ instead of $d^2(W_0^2-L_0^2G)^2$ saves one more. Finally
\[
 u_r^2+\tau u_rv_r+\kappa v_r^2
 =u_r(u_r+\tau v_r)+\kappa v_r^2
\]
saves a third multiplication. These identities hold over every commutative ring and every coefficient specialization. Thus the source's entire root-recovery proof, including its auxiliary rational Pell saturation and weak-unit analysis, transfers without modification~\cite{positive30}.

For an additional one-operation saving in the system, use $z_{\rm old}=z+1$. Transport becomes
\[
 (K_0+w)C=F+z(q-1).
\]
This gives a 103-operation system with three equations, maximum residual degree 1072. Its positive zero set is in bijection with the old one: the inherited recovery gives $C>0$, $X=wq=2^R$, $R>3q$, and $F<q$. Hence $w>q>F$ and
\[
 (z_{\rm old}-1)(q-1)=(K_0+w)C-F>0,
\]
so $z_{\rm old}\ge2$. The shift is applied to this system only; it is not silently inserted into the 108-operation product.

The new tuple below changes the actual fixed-coefficient polynomial when compared with package 30's original tuple. The three circuit rewrites are identities with ports held equal; semantic rebinding is a separate complete-tuple theorem. These two statements should not be conflated.

\section{Smaller materialized coefficients}\label{sec:coeff}
\subsection{Geometry and the finite recipe}
The new complete native tuple is
\[
 \beta=B-1,\qquad K_0=D_C+BD_R,\qquad M_C,
 \qquad M_F=M_{F,\mathrm{native}}+B-1,
\]
where
\begin{equation}\label{eq:newbase}
 \rho=2^{121},\qquad L_c=8801,\qquad
 d=121\cdot8801=1064921,\qquad B=2^d.
\end{equation}
The literal 943-tile language is regenerated from the 430 transitions of K215 and checked against the supplied ordered tile list. Its mixed-field layout uses clause bands $T_1=4393$, $T_2=5495$, and spatial shift 42. The arithmetic generator and \smallpath{compiler/COEFFICIENT_PROOF.md} provide the full coefficient dictionary and binary masks; no unspecified optimization or search is part of their definition.

The selector fields, bit-copy fields, two clean clause bands, and synchronization anchors are the mixed layout from package 25. The uniform convolution mass bound has 118 bits, and its computational carry threshold has 120 bits, so the odd radix width 121 suffices. The largest computation-support exponent is 8791. Put 860 extra mask bits in zero cells 8792 through 8799: seven fields of $2^{119}-1$, followed by one field of $2^{27}-1$. Their separate margin is $2(2^{119}-1)+4<2^{121}$. This preserves the computation support and balances the native population at 989.

The independent checker verifies 69,230 basis convolution conditions, the disjoint clause bands, the unique anchor intersection, the carry margin, and
\begin{gather*}
 0<M_C,M_{F,\mathrm{native}}<B-1,\qquad
 M_C\equiv2\pmod4,\quad M_{F,\mathrm{native}}\equiv4\pmod8,\\
 \operatorname{popcount}(M_C)+\operatorname{popcount}(M_{F,\mathrm{native}})=d.
\end{gather*}
These checks discharge the finite layout hypotheses. They alone do not establish completeness; the next argument replaces the old power-of-five padding assumption.

\subsection{A complete Boolean padding controller}
Factor $d=11^2\cdot13\cdot677$ and choose the prime $p=2341$ and step $k_c=2340=p-1$. Exact modular arithmetic gives
\[
 \operatorname{ord}_d(B)=780\mid2340,\qquad
 B^{2340}\bmod2341^2=2556373\ne1,
\]
with the latter residue equal to one modulo 2341. Therefore
$v_{2341}(B^{2340}-1)=1$. Define $H_0=1+BD_C+B^2(D_R+1)$. The required unit residues are:
\begin{center}
\begin{tabular}{rrrr}
\toprule
Prime & $B\bmod p$ & $(B^2-1)\bmod p$ & $H_0\bmod p$\\
\midrule
11&2&3&9\\13&6&9&7\\677&611&293&654\\2341&2100&1896&1225\\
\bottomrule
\end{tabular}
\end{center}
Choose sufficiently large independently padded rectangular tableau dimensions $w_c,h_c$, both powers of 2341, and $N_c=w_ch_c=2341^n$. Require $N_c>d\,2341^2$ and $h_c>2341$. Because every relevant order divides 2340,
\[
 q=B^{N_c}\equiv B^{w_c}\equiv B\pmod{d\,2341}.
\]
Switching the lower dummy bit in cell $j$ changes the actual packed index by
\[
 \delta R=-\Gamma B^j,\qquad
 \Gamma=\rho^{1063}(q^2-1)
 [1+q(D_C+BD_R+B^{w_c})].
\]
The residue table makes $\Gamma$ a unit modulo $dN_c$.

For completeness, every residue modulo $dN_c$ is represented by a Boolean subset of
\[
 \{B^{2+2340j}:0\le j<T\}\ \cup\ \{B^3\},\qquad
 T=N_c/2341.
\]
To prove this, put $C_0=d\,2341$, so $dN_c=C_0T$ and $T>C_0$. The first $T$ powers of $B^{2340}$ are exactly the $T$ residue classes $1+C_0z$ modulo $C_0T$, by the valuation and order calculation. Normalize the target by $B^{-2}$ and optionally subtract $B$ so that it is nonzero modulo 2341. Write its residue modulo $C_0$ as $a$, with $1\le a<C_0<T$ and $\gcd(a,T)=1$, and write the target as $a+C_0u$. Solve
\[
 az_0+\frac{a(a-1)}2\equiv u\pmod T.
\]
The $a$ distinct additive orbit coordinates $z_0,z_0+1,\ldots,z_0+a-1$ have the required sum. Include $B^3$ if the earlier subtraction was made. No Boolean position is reused.

The largest available cell is $N_c-N_c/2341-2341+3$. Since $w_c<N_c/2341$, its required spatial and temporal translates do not wrap. Thus the chosen dummy bits can enforce
$R\equiv dw_c\pmod{dN_c}$ without altering a tile, overlap, or marker. This restores the actual rotation required by $X=2^R$. The inherited carry margin preserves positive transport slack, and the unchanged Pell converse completes the positive witness. This is package 24's arbitrary-factor controller specialized to the smaller mixed geometry, with every numerical condition checked anew.

\subsection{Payloads and portability}
\begin{center}
\begin{tabular}{lrrr}
\toprule
Port & Bit length & Raw bytes & Gzip bytes\\
\midrule
$\beta$&1,064,921&133,116&164\\
$K_0$&1,070,004&133,751&9,794\\
$M_C$&1,064,921&133,116&264\\
$M_F$&1,064,922&133,116&263\\
\midrule
Total&&533,099&10,485\\
\bottomrule
\end{tabular}
\end{center}
The former smallest displayed base exponent was $5^9=1,953,125$. The new exponent is about 45.48 percent smaller. The four uncompressed payloads fall from 977,220 to 533,099 bytes, a 45.45 percent reduction. Compression size is a storage measure, not coefficient height. All four ports must be replaced together, and $M_F$ is already shifted.

The common native interface, including its completeness proof above, allows this tuple to bind the twelve-positive-variable forms and inherited direct positive baselines. The integer Jones families in this archive retain their T4 recipe. We make no unverified claim that the new tuple has been substituted into their fixed quartic/index construction.

\section{Verification, reproducibility, and limitations}\label{sec:verification}
\subsection{Exact artifacts}
The archive includes the complete LaTeX article and compiled PDF; literal division-free circuits for the fifteen- and sixteen-integer equations and the twelve-positive-variable product/systems; the complete coefficient recipes and four compressed materialized small ports; source provenance and licences; and standard-library Python builders and checkers. \texttt{verify\_all.py} runs the delivered verification in a fresh extraction. No repository checkout, credential, or network connection is required.

The integer evaluator computes the same Jones code through a factored digit-polynomial algorithm. It includes the background $z_*$ contribution from every degree-at-most-four exponent vector, not only the finitely listed nonzero coefficients. Shared fixed powers, degree-four generating recurrences, coefficient aliases, and a factored sparse evaluation avoid repeatedly constructing large powers. This accounts for most of the reduction from the predecessor's 20,607 operations; the new mathematics accounts for the removed guards and smaller weight. Both kinds of improvement are credited separately.

\subsection{Degree certificates}
For the sixteen-variable body, with $L=5^{33}$, the inherited exact degrees are
\[
 \deg m=2221L+11,\quad\deg a_0=3800L+14,\quad
 \deg a_1=24640L+68,\quad\deg D_0=6875L+20.
\]
The first five guard products have recurrence $c_j=3c_{j-1}+3$, giving $c_5=363$. Therefore
\begin{center}
\begin{tabular}{lr}
\toprule
Expression & Exact degree\\
\midrule
$C_*$ and $w_1$ & $2221L+1102$\\
$w_2$ & $4121L+1110$\\
$S$ & $34375L+107$\\
$T$ & $34375L+105$\\
$X$ & $34375L+108$\\
$A_2$ & $101632L+2502$\\
\bottomrule
\end{tabular}
\end{center}
$A_2$ has strictly greater degree than $X^2,A_0,A_1$. Hence the output leading term is $A_2^4$, giving $406528L+10008$. The checker tracks a nonzero diagonal leading coefficient modulo 1,000,003 independently of every unspecified coefficient leaf. This is an exact certificate of a nonzero leading homogeneous component, not a probabilistic assertion of degree.

For the fifteen-variable circuit, the new radix has degree two and a positive fixed leading coefficient $b_*$. Its separate certificate records the dependence on $b_*$ explicitly and proves exact degree \DegreeFifteenFormula. It does not substitute a zero value for that fixed positive coefficient or treat a generic degree bound as automatically attained. All circuit nodes are checked for a valid topological order, live variables, and the stated operation ledger.

\subsection{What the finite checks establish}
The positive rewrites are proved as exact sparse polynomial identities, with additional signed modular circuit comparisons. The integer checks compare the optimized body to the literal source, evaluate the finalizer independently as an eight-factor product, check actual coefficient residues, and audit the degree ledger. Small-modulus intersectivity and Pell-orbit examples exercise implementation paths; the unbounded claims rest on the proofs in this report. The coefficient audit independently reconstructs the payloads without importing the builder and checks their hashes, support geometry, population, controller prime, orders, and unit residues.

No checker searches for a ZFC contradiction or certifies the actual absence of a root. No new whole-project Lean/Rocq build is claimed. The end-to-end provability assertion uses the inherited compiler and calibration as mathematical theorems; this is the same proof boundary explicitly recorded by the source collection. The archive makes that boundary reviewable rather than concealing it behind a numerical PASS label.

\section{Further research questions}
\begin{enumerate}[leftmargin=*,itemsep=6pt]
\item \textbf{Formalize cofinal radix completeness.} Generalize the source \texttt{USys\_of\_mem} statement to accept any power of two exceeding the explicit finite bounds. This is a small, precise interface improvement and directly supports the fifteen-variable theorem.
\item \textbf{Formalize the Pell orbit guard removal.} Lemma~\ref{lem:orbit}, pairwise coprimality, and the $j$ replacement can be proved independently of the compiler. Joining these to the existing ratio theorem would substantially reduce the unformalized part of the signed construction.
\item \textbf{Eliminate further sign guards only with a signed soundness theorem.} The surviving guards on $g,s,w,\phi$ serve different purposes. A positivity argument for one must preserve the ratio margins and congruence classes; simply replacing a variable by its absolute value or square need not preserve witnesses.
\item \textbf{Combine the integer reduction with the 31-variable quartic.} Rebuilding the Jones index at $\nu=31$ should lower the dominant degree by roughly a factor of five, but it requires a fresh coefficient/index binding and proof of every specialization condition. This archive deliberately retains the authenticated T4-bound $\nu=32$ source.
\item \textbf{Instantiate the ten-integer construction.} Sun's theorem gives a lower abstract bound. A full explicit source-code-to-polynomial specialization for the calibrated refutation predicate would provide a better concrete integer count, with an independently reviewable coefficient recipe.
\item \textbf{Optimize geometry and evaluation separately.} The new coefficient height comes from geometry and padding, while the 108-operation positive polynomial comes from circuit identities. Further work should report both; a reduction in compressed file size alone does not simplify the mathematical coefficient height.
\item \textbf{Test the remaining frontiers.} A positive quartic with fewer than 31 unknowns, a direct positive equation below 67 operations, or a proof of a meaningful lower bound in a restricted circuit model would each improve a different part of the collection. No such improvement is asserted here.
\end{enumerate}

\section*{Attribution and source policy}
This report and its new arguments were prepared with OpenAI assistance for Vladimir Reshetnikov's ProveIt project. Source constructions are attributed by package and immutable commit; no attribution of human authorship or independent external review is implied. The new material is offered under the repository's MIT-0 convention. Copied machine and source files retain their supplied notices and licences. The public repository was read, not modified; the package is prepared for incorporation.

\begingroup\small\raggedright
\begin{thebibliography}{20}
\bibitem{collection} V.~Reshetnikov, \emph{ProveIt: Con(ZFC) in Diophantine form}, reconciliation note and packages 24--30, commit \pin. Repository path \smallpath{SetTheory/Cardinals/docs/reports/hilbert-tenth-problem/con-zfc-diophantine-forms/}. \url{https://github.com/VladimirReshetnikov/ProveIt}.
\bibitem{signed17} ProveIt package 29, \emph{Seventeen unrestricted integers: complete new reduction}, \smallpath{29-simpler-forms/signed17/SIGNED17_PROOF.md}, at the comparison commit. Original T4/Jones source files and hashes are preserved in the archive.
\bibitem{positive30} ProveIt package 30, \emph{A smaller twelve-positive-variable synthesis}, \smallpath{30-quartic31/positive/POSITIVE12_PROOF.md}; and \emph{A 31-positive-variable quartic by eliminating a packing register}, \smallpath{30-quartic31/arithmetic/QUARTIC31_PROOF.md}.
\bibitem{mixed25} ProveIt package 25, \emph{Smaller Diophantine forms of ZFC consistency}, mixed125 coefficient proof and binding, \smallpath{25-hybrid12/coefficients/}.
\bibitem{padding24} ProveIt package 24, \emph{Smaller explicit Diophantine forms of ZFC consistency}, coefficient section and arbitrary-factor padding argument, \smallpath{24-integer24/sections/coefficients.tex}.
\bibitem{fast27} ProveIt package 27, \emph{Smaller Diophantine presentations of ZFC consistency}, optimized nine-variable evaluator and proof, \smallpath{27-wave4/nine/}.
\bibitem{ninelean} ProveIt, \emph{Jones 1982 formalization}, \smallpath{Computability/HilbertTenthProblem/Lean/Diophantine/Paper1982/}, especially \texttt{Master1.lean}, \texttt{Master3.lean}, \texttt{NineReduce.lean}, \texttt{NineRatio.lean}, and \texttt{NineMain.lean}, at the comparison commit.
\bibitem{endpoint} ProveIt, \texttt{KEndpoint215.blank\_halts\_iff\_zfc\_refutation}, \smallpath{Computability/BusyBeaver/Lean/BusyBeaver/ZFC/KEndpoint215.lean}; literal machine \smallpath{Computability/BusyBeaver/Research/ZFC/candidate-k215/k215.tm}.
\bibitem{jones} J.~P.~Jones, Universal Diophantine equation, \emph{Journal of Symbolic Logic} \textbf{47} (1982), 549--571. The explicit instance used here is supplied and audited by the pinned ProveIt source.
\bibitem{sun} Z.-W.~Sun, \emph{Ten unknowns for Hilbert's tenth problem over the integers}, arXiv:2609.01594v2, 2 September 2026. \url{https://arxiv.org/abs/2609.01594v2}. Cited for the general ten-integer representation theorem and related relation-combining context, not as an instantiated coefficient construction in this archive.
\end{thebibliography}
\endgroup
\end{document}
